% from ega4-17-fr.tex
\oldpage[IV-4]{84}

\emph{a fortiori} rationnel sur $k$, et comme $\sh{O}_{Y,y}$ est alors une $k$-algèbre formellement lisse pour les topologies discrètes et que $\sh{O}_{Y,y}/\mathfrak m_y=k$, il résulte de (0, 20.5.14) que l'homomorphisme canonique $\mathfrak m_y/\mathfrak m_y^2\to(\Omega^1_{Y/S})_y\otimes_{\sh{O}_{Y,y}}k(y)$ est bijectif; de même, l'homomorphisme canonique $\mathfrak m_x/\mathfrak m_x^2\to(\Omega^1_{X/S})_x\otimes_{\sh{O}_{X,x}}k(x)$ est bijectif. L'hypothèse c'), équivalente à c), signifie donc ici que l'homomorphisme canonique
\[
(\mathfrak m_y/\mathfrak m_y^2)\otimes_{k(y)}k(x)\to\mathfrak m_x/\mathfrak m_x^2
\]
est injectif. Comme l'anneau $\sh{O}_{X,x}$ est régulier, la conclusion résulte de (0, 17.3.3). C.Q.F.D.

\begin{corollary}[17.11.2]
\label{IV.17.11.2}
Sous les hypothèses générales de (17.11.1), les conditions suivantes sont équivalentes:
\begin{enumerate}
  \item[a)] $f$ est lisse au point $y$ et $g$ est étale au point $x$.
  \item[b)] $h$ est lisse au point $x$ et $g$ est étale au point $x$.
  \item[c)] $h$ est lisse au point $x$, et l'homomorphisme canonique
  \[
  (g^*(\Omega^1_{Y/S}))_x\to(\Omega^1_{X/S})_x
  \]
  est bijectif.
  \item[c')] $h$ est lisse au point $x$, et l'homomorphisme canonique
  \[
  ((\Omega^1_{Y/S})_y\otimes_{\sh{O}_{Y,y}}k(y))\otimes_{k(y)}k(x)\to(\Omega^1_{X/S})_x\otimes_{\sh{O}_{X,x}}k(x)
  \]
  est bijectif.
\end{enumerate}
Supposons en outre que l'homomorphisme $k(y)\to k(x)$ soit bijectif. Alors les conditions précédentes sont aussi équivalentes à la suivante:
\begin{enumerate}
  \item[d)] $h$ est lisse au point $x$, et l'application canonique $\mathrm{T}_{X/S}(x)\to\mathrm{T}_{Y/S}(y)$ (16.5.12) est bijective.
\end{enumerate}
\end{corollary}

Chacune des conditions a), b), c) de (17.11.2) est équivalente à la conjonction de la condition correspondante de (17.11.1) et du fait que $g$ est non ramifié au point $x$, compte tenu de (17.2.2) en ce qui concerne la condition c); d'où l'équivalence de a), b) et c). L'équivalence de c) et c') résulte de l'équivalence des conditions correspondantes de (17.11.1) et du lemme de Nakayama; l'équivalence de c') et d) lorsque $k(x)=k(y)$ est immédiate par transposition (16.5.12).

\begin{corollary}[17.11.3]
\label{IV.17.11.3}
Soient $h:X\to S$ un morphisme lisse, $s_i$ ($1\leq i\leq r$) des sections de $\sh{O}_X$ au-dessus de $X$ (qui sont aussi des sections de $h_*(\sh{O}_X)$ au-dessus de $S$), $g:X\to S[T_1,\ldots,T_r]=\mathbf{V}_S^r$ le $S$-morphisme correspondant à l'homomorphisme de $\sh{O}_S$-Modules $\sh{O}_S^r\to h_*(\sh{O}_X)$ défini par ces sections (II, 1.2.7). Pour que $g$ soit lisse (resp. étale) en un point $x\in X$ il faut et il suffit que les $(d_{X/S}(s_i))_x$ fassent partie d'une base (resp. constituent une base) du $\sh{O}_{X,x}$-module $(\Omega^1_{X/S})_x$.
\end{corollary}

Il suffit d'appliquer (17.11.1) (resp. (17.11.2)) en prenant pour $f$ le morphisme structural $S[T_1,\ldots,T_r]\to S$.

\begin{corollary}[17.11.4]
\label{IV.17.11.4}
Pour qu'un morphisme $h:X\to S$ soit lisse en un point $x\in X$, il faut et il suffit qu'il existe un voisinage ouvert $U$ de $x$, un entier $r$, et un $S$-morphisme étale $U\to S[T_1,\ldots,T_r]$.
\end{corollary}

